\section*{Introduction}

Why are strict religious groups often unusually cohesive and durable? The canonical answer begins with religious organizations as clubs that produce collective goods. In Iannaccone's account, apparently wasteful demands can make these clubs stronger: costly requirements involving dress, diet, ritual, time, and money screen out people whose low commitment would otherwise permit them to free-ride on the sacrifices of others \citep{iannaccone1988,iannaccone1992,iannaccone1994}. Those who remain redirect resources away from competing activities and toward the group, increasing both the average commitment of members and the value of membership itself. Strictness, on this view, is not organizational pathology but a solution to adverse selection and weak participation. This club-goods framework has become the dominant economic explanation of religious strictness, as reviews of the economics of religion emphasize \citep{iyer2016,carvalho2019}. It has also generated an unusually broad family of applications. Sacrifice and nontransferable religious capital help explain ultra-Orthodox Jewish education and fertility \citep{berman2000}; monitoring institutions can sustain religious production \citep{mcbride2007}; costly signals can resolve commitment problems in religious militancy \citep{bermanlaitin2008}; and veiling can operate as a commitment device that changes social interaction and marriage-market incentives \citep{carvalho2013}. Across these applications, costly observances select committed participants and align their investments with the group.

Yet the same arrangements have another face. Every mechanism that screens at entry can simultaneously raise the cost of exit after entry. A social network concentrated inside the group is valuable to a committed member, but it is also difficult to replace after departure. An identity made inseparable from membership can deepen solidarity, but it can also make departure psychologically and socially disorganizing. Distinctive dress, diet, language, schooling, and ritual can certify commitment among insiders while isolating members from outsiders who might provide employment, friendship, marriage, or recognition. Screening and exit foreclosure are therefore the same social structure described from opposite directions. Looking inward from the organizational boundary, strictness selects and coordinates participants; looking outward from the position of a member considering departure, it degrades alternatives and makes separation costly. Case scholarship often recognizes both processes at once without analytically separating them. An analysis of the Skoptsy sect, for example, argues that high entry costs screened for devoted individuals while high exit costs reduced the viable alternatives for living outside the sect [VERIFY: citation key for the Skoptsy sect analysis]. The two mechanisms can coexist, but coexistence does not establish which one accounts for durability.

That ambiguity is now a live theoretical problem rather than a terminological distinction. \citet{bok2025} observes that research on religious exit has concentrated mainly on who leaves, why they leave, and what follows, while the forms and roles of exit costs have received comparatively little systematic theorization. The systematic theorization of the exit-cost side is therefore recent relative to the mature literature on the determinants and consequences of exit. The field has a mature screening account and a newly articulated exit-cost problem, but no design that can adjudicate their relative causal force. The identification problem is plain. In observational data, screening and exit foreclosure arrive bundled together. Any strict group that persists has already selected some people willing to bear its demands and has also organized relationships, identities, and opportunities in ways that make leaving consequential. Comparing strict and lenient groups cannot recover the separate contribution of either process, because the strictness that filters entrants also transforms the outside option of those who stay. Nor can studying former members solve the problem by itself: observed leavers are selected on having overcome the very costs whose preventive effect is at issue. Without independent variation, evidence consistent with successful screening is also consistent with blocked exit.

Agent-based simulation supplies the missing instrument because it permits the two structures to be varied as distinct inputs while measuring any coupling between their effects. An experimental design can hold the observability and rewards that support enforcement constant while changing members' capacity to leave, or hold exit conditions constant while changing whether an enforcement apparatus can activate. No observational design can cleanly impose either intervention on otherwise equivalent religious communities. The strategy follows an established computational-religion program associated with Shults and Wildman [VERIFY: publication year and citation key for Shults and Wildman], formal arguments for bringing simulation into the study of religion \citep{wildmansosis2011}, computational work linking cognition, culture, and social dynamics \citep{gorelemosshultswildman2018}, and agent-based analysis in the economics of religion \citep{iannacconemakowsky2007}. Simulation here is therefore not a substitute for historical or ethnographic evidence. It is a method for identifying structural effects that those forms of evidence necessarily observe in combination.

This paper distinguishes two structural dimensions that are largely separable, but coupled at the frontier margin: whether an enforcement apparatus activates, and whether members can leave. The quantitative caveat to this separation is reported in the results. Three findings result. First, activation follows a joint threshold in code observability and the reward for enforcement. Neither condition is sufficient alone. Across a 63-cell grid crossing nine observability levels with seven reward levels, enforcement appears along a frontier defined jointly by both parameters rather than along either axis independently. Second, exit closure is necessary but not sufficient for capture. In the exogenous closure experiment, the capture rate rises monotonically from 0 at fully open exit to 0.889 at the maximum imposed closure. Yet the low-observability, low-reward cell remains at zero capture at every closure level, including the maximum. Closing exit can retain an apparatus that activates, but it cannot create one where enforcement lacks both legibility and reward. Third, the concentration of punishment that appears to mark enforcement-dominant regimes is largely manufactured by institutional privilege. In the floor-versus-ceiling ablation, the top 5 percent of active agents issue 0.091 of punishments without the privilege bundle and 0.580 with it. Relative to the ceiling, 0.843, or roughly 84 percent, of top-5-percent concentration is privilege-manufactured. The floor Gini is 0.240, compared with 0.028 under a same-volume random-allocation null. These results distinguish the conditions that form enforcement, retain its subjects, and concentrate coercive activity. Exit capacity is treated as an exogenous structural input as a methodological safeguard. A specification in which outside-option degradation was driven by enforcer share was circular because enforcer share entered the capture criterion. Once that coupling was removed, capture occurred nowhere in the tested space, motivating the exogenous treatment used for identification.

The concentration result also qualifies a prominent interpretation in the punishment literature. The costly-punishment tradition asks how sanctioning can arise through decentralized interaction even when punishers bear private costs. Experimental and evolutionary models show how punishment can sustain cooperation and how sanctioning behavior can spread or persist under specified population and institutional conditions \citep{fehrgachter2000,boydgintisbowles2010,hauertetal2007}. From that perspective, a skewed distribution of punishment may look like an emergent product of heterogeneous strategies, repeated interaction, or selection among decentralized actors. The ablation suggests a different interpretation for delegated architectures. When one role enjoys exclusive or strongly asymmetric authority, lower costs and risks, and material reinforcement, concentrated sanctioning is substantially built into the opportunity structure before interaction generates any further inequality. This is not a general claim that punishment concentration never emerges from decentralized dynamics. It is a bounded contrast: in models with an institutionally privileged enforcement role, observed concentration cannot be treated as evidence of decentralized emergence without first removing or measuring the privileges that allocate coercive opportunities.

The mechanism is content-neutral. It is defined over structural properties: how observable a code's requirements are, what rewards attach to enforcing them, whether sanctioning authority is delegated asymmetrically, and how costly it is for a member to leave. Nothing in this definition depends on the truth, theology, or substantive moral content of a doctrine. The same logic can apply to any organization that codifies conduct and attaches status to compliance, including political movements, residential communities, firms, schools, or ideological associations. Religious regimes provide the empirical terrain because they offer the longest and richest documented record of codified conduct, costly membership, delegated moral authority, and contested exit. They also make the duality especially visible, since practices experienced as identity, discipline, and mutual commitment by some members may function as barriers to outside life for others. The paper makes no claim that codified belief tends toward coercion. Its claim is conditional and structural: where observable rules, enforcement rewards, delegated privileges, and constrained exit combine, they generate dynamics that cannot be inferred from doctrine alone.

The remainder of the paper develops the argument in four steps. The next section connects screening, exit costs, and costly punishment and derives the hypotheses. The methods section specifies the agent-based design and the independent interventions on activation, exit, and privilege. The results section presents the activation frontier, exit-closure dose-response, and privilege ablation. The conclusion states the empirical implications, limits, and tests that follow from separating selection into a group from members' capacity to leave it.
